% !TeX root = main.tex
\chapterimage{chapter-01-introduzione.jpg}
\chapter{Introduzione}\label{capitolo-1-introduzione}

Insegnante che ti appresti a leggere questa guida con curiosità, scetticismo, speranza, timore, entusiasmo o un mix di tutte queste cose: un sorriso a te! \emoji{smiling-face-with-smiling-eyes}

Questa guida nasce originariamente per proporre attività di \textbf{informatica} per la \textit{scuola primaria}, ma le proposte sono utilizzabili (o facilmente adattabili) anche per la \textit{scuola secondaria di primo grado}, soprattutto se le studentesse e gli studenti non hanno mai affrontato prima questi argomenti.

Leggendola, a volte ti verrà spontaneo pensare: ``ma non è così diverso da quello che faccio già!'': hai perfettamente ragione! Ciò che faremo in questo percorso insieme è rileggere queste attività  con una nuova lente, quella dell'informatica, che potremmo chiamare \emph{computazionale}.

Molte attività si basano infatti su esperienze motorie, sull'esplorazione progressiva dello spazio ambientale e personale, sull'uso intenzionale del linguaggio, sulla collaborazione tra pari e sulla riflessione sugli errori: tutte pratiche che, adattate all'età e all'esperienza della classe, accompagnano e sostengono l'apprendimento nel primo ciclo. La novità sta nella prospettiva con cui tali pratiche vengono affrontate, ossia attraverso la lente dell'informatica. Inoltre, le stesse attività diventano occasioni per rafforzare le competenze di risoluzione dei problemi, sviluppare la capacità di pianificazione e, al tempo stesso, potenziare la creatività di alunni e alunne.

Questo è il punto di partenza di un percorso da affrontare per tutto il primo ciclo. Prima, attività senza computer, semplici e giocose; poi, l'introduzione graduale di strumenti digitali facili da usare e gratuitamente disponibili in rete. Nella scuola primaria questi strumenti possono essere introdotti a partire dalla terza; nella secondaria la progressione va adattata alle esperienze precedenti della classe.

\section{Perché l'informatica fin dalla scuola primaria }\label{perchuxe9-linformatica-fin-dalla-scuola-primaria}

\subsection{Società digitale }\label{societuxe0-digitale}

Oggi viviamo in un mondo in cui le tecnologie digitali plasmano il modo in cui ci informiamo, ci intratteniamo, comunichiamo con gli altri e prendiamo decisioni nella vita quotidiana. Per orientarsi in questa realtà in modo consapevole non basta saper usare tablet, computer o applicazioni: è necessario capirne i meccanismi sottostanti, così da poterli dominare e valutare criticamente.

La diffusione dell'intelligenza artificiale rende questa esigenza ancora più urgente: sempre più spesso affideremo alle macchine compiti complessi, da svolgere seguendo le nostre indicazioni. Per questo motivo è fondamentale imparare a interagire con la tecnologia digitale non solo come consumatori passivi, ma come protagonisti attivi e consapevoli.

Qui entra in gioco l'informatica. È bene chiarire subito un equivoco diffuso: l'informatica non coincide con la tecnologia. Fare informatica non significa semplicemente usare dispositivi digitali, ma studiare i principi che li rendono possibili. È una disciplina scientifica a tutti gli effetti, che offre ai nostri alunni e alunne un nuovo modo di pensare e di interpretare il mondo. Così come la lettura apre le porte al linguaggio dei testi e la matematica permette di comprendere quantità e relazioni, l'informatica consente di padroneggiare i principi e i linguaggi alla base delle tecnologie digitali che ci circondano.

In quest'ottica--oltre al raggiungimento degli obiettivi disciplinari, così come accade per le altre materie--le attività proposte mirano a sviluppare competenze per:

\begin{itemize}
\item
  \textbf{Affrontare problemi} scomponendoli in passi chiari e risolvibili;
\item
  \textbf{Descrivere soluzioni} che possano essere eseguite da altri;
\item
  \textbf{Riflettere sul proprio ragionamento}, allenando metacognizione e pensiero critico;
\item
  \textbf{Usare la creatività} per immaginare soluzioni nuove e originali.
\end{itemize}

Avviare l'alfabetizzazione informatica fin dalla scuola primaria significa piantare le fondamenta di un curriculum verticale che accompagni alunni e alunne per tutto il loro percorso scolastico. Il primo ciclo è infatti il luogo in cui tutti, indipendentemente dal contesto familiare o dalle opportunità offerte dall'esterno, possono accedere alle stesse esperienze di apprendimento. Introdurre l'informatica in questo contesto significa iniziare un percorso di equità, offrendo a ciascuno le basi per un rapporto attivo e critico con la tecnologia, accrescendo la consapevolezza dei meccanismi che regolano gli strumenti quotidiani e contribuendo a ridurre il divario digitale.

\subsection{Informatica nelle Indicazioni Nazionali 2025 }\label{informatica-nelle-indicazioni-nazionali-2025}

Le Indicazioni Nazionali 2025 \parencite{mim-indicazioni-2025} introducono, per il primo ciclo, competenze, obiettivi e conoscenze di Informatica nelle discipline Matematica e Tecnologia.

Le Indicazioni Nazionali riconoscono l'importanza di avvicinare fin da subito alunne e alunni ai concetti fondamentali della disciplina.

Nella parte di matematica, troviamo obiettivi relativi a tre aree fondamentali: informazioni, dati e loro rappresentazione; algoritmi; programmazione. Queste tre aree, che riguardano l'informatica, sono coperte da questa guida. Verranno discusse in dettaglio e proposte attività didattiche relative alle competenze, agli obiettivi e alle conoscenze presenti nelle Indicazioni Nazionali, specificati alla fine di ogni capitolo\footnote{I riferimenti curricolari riportati alla fine dei capitoli riguardano la scuola primaria, da cui nasce il volume; l'adattamento delle attività alla secondaria richiede di calibrarne consegne e approfondimenti sul livello della classe e sulle esperienze pregresse degli alunni e delle alunne.}.


Nella parte di tecnologia, troviamo obiettivi relativi a varie aree fondamentali: aspetti tecnologici (componenti hardware e software di un sistema informatico, rete Internet e i suoi servizi); aspetti di sicurezza (protezione dei dati personali, crittografia) e uso consapevole delle tecnologie; creatività digitale, specialmente tramite la realizzazione di semplici applicazioni informatiche. In questa guida ci concentreremo in particolare sull'ultima (tramite la programmazione creativa con Scratch) e faremo accenni alle altre, rinviando gli approfondimenti a fasi successive del percorso.

Per una lettura ragionata dei contenuti informatici delle Indicazioni Nazionali 2025, rimandiamo a \textcite{lodi-lonati-2026}.

\section{Perché questo manuale }\label{perchuxe9-questo-manuale}

Portare l'informatica in classe significa introdurre un linguaggio e un modo di pensare che non fanno parte dell'esperienza scolastica o universitaria di molti insegnanti. In questo territorio può capitare di sentirsi senza mappe né punti di riferimento chiari.

Questo manuale nasce per colmare proprio quel vuoto. Non solo per offrire spunti per attività da far svolgere in classe e a casa, ma anche per dare un filo conduttore, una visione unitaria e coerente di cosa significhi fare informatica nel primo ciclo.

Troverai:

\begin{itemize}
\item
  proposte di attività da svolgere in classe, adatte all'età e integrate nella vita di classe, e strategie didattiche per metterle in pratica al meglio;
\item
  consigli pratici, soluzioni ed estensioni per le attività proposte;
\item
  la spiegazione, rigorosa ma il più possibile accessibile, dei concetti fondamentali di informatica necessari per capire e lavorare su ciò che è richiesto nelle Indicazioni Nazionali 2025.
\end{itemize}

La nostra ambizione è duplice: fornire strumenti affidabili con cui lavorare e, allo stesso tempo, incoraggiare la sperimentazione e l'adattamento. Perché l'informatica nel primo ciclo non è solo una nuova disciplina da introdurre: è un'opportunità per aprire spazi di scoperta, di creatività e di ragionamento condiviso, in cui anche l'insegnante può imparare insieme ad alunni e alunne.

\section{Dalla teoria alla pratica: leggi qui!}

Fermiamoci un attimo! Sfogliando le prossime pagine, potresti avere la sensazione di trovarti davanti a tanto, tantissimo testo. \emoji{worried-face} \emoji{neutral-face}

Scusaci: in questi primi due capitoli ci prendiamo un po' di spazio per costruire un linguaggio comune, chiarire i concetti fondamentali e collocarli nel quadro educativo del primo ciclo. Non temere: non sarà così per tutto il manuale.

A partire dal Capitolo~\ref{algoritmi-quotidiani}, proponiamo attività pratiche da svolgere in classe. Se vuoi, puoi cominciare a dare un'occhiata direttamente alle attività: troverai tanti termini, alcuni familiari, altri nuovi, altri ancora usati anche in altre discipline, ma che in informatica assumono una sfumatura diversa. Dare un'occhiata alle attività ti spingerà a tornare indietro e a leggere con maggiore attenzione e consapevolezza queste pagine introduttive.

I box contrassegnati da questi simboli ti aiuteranno durante le attività:

\begin{itemize}
\item[\guideiconAttention{}]
  attenzioni particolari per valorizzare l'esperienza e guidare l'osservazione;
\item[\guideiconQuestion{}]
  domande stimolo, per trasformare i giochi in un'occasione di riflessione informatica.
\end{itemize}

Il nostro obiettivo è semplice: offrirti una raccolta di strumenti da usare. La parte scritta in questo manuale serve a darti sicurezza e una cornice di senso; il cuore del percorso sono le attività pratiche, in cui potrai sperimentare insieme ad alunni e alunne l'informatica come disciplina, mettendo in pratica con essa un uso del linguaggio creativo, rigoroso e divertente.

\subsection{Informatica, pensiero computazionale e coding: facciamo chiarezza }\label{informatica-pensiero-computazionale-e-coding-facciamo-chiarezza}

Sei ancora qui! Molto bene \emoji{smiling-face-with-smiling-eyes}

Per dare coerenza a questo percorso serve prima di tutto una bussola. A scuola, infatti, raramente si parla di informatica, mentre capita spesso di sentire parlare di pensiero computazionale, programmazione e \textit{coding} come se fossero sinonimi. Chiarire il significato di questi termini è il primo passo per dare all'informatica il giusto spazio e valore, evitando confusioni e fraintendimenti.

Questo manuale parla di informatica: la disciplina scientifica che studia l'\textbf{elaborazione automatica dell'informazione}--una disciplina autonoma con idee, metodi e un linguaggio propri. Come dici? Pensavi che l'informatica studiasse i computer? Questa definizione ti fa pensare a un'intelligenza artificiale che sceglie autonomamente quali informazioni dare al telegiornale? Non è proprio così. Pensa all'informazione come a qualunque cosa utile (per esempio, la foto dell'ultimo tuo viaggio è utile per fare invidia agli amici) che vuoi far elaborare in modo automatico (per esempio applicando un filtro del tuo smartphone che fa sembrare il cielo azzurro anche se in realtà era nuvoloso). Nel capitolo successivo, proveremo a dare definizioni più chiare e precise.

Così come in matematica si parla di ``pensiero matematico''--caratterizzato da modelli, astrazioni e strategie tipiche del ``fare matematica''--anche l'informatica permette di sviluppare un modo di ragionare disciplinare, che potremmo chiamare ``pensiero informatico''. Questo modo di pensare è caratterizzato dalla capacità di affrontare una classe di problemi, individuare una soluzione generale e descriverla in modo così chiaro e preciso che questa può essere eseguita da un agente esterno, spesso con un insieme molto preciso e limitato di azioni disponibili. In queste pagine useremo però il termine più diffuso, \textbf{pensiero computazionale}: non si tratta di ``pensare come un computer'', ma di adottare un metodo di ragionamento tipico dell'informatica, da applicare alla soluzione di problemi che ci troviamo davanti--a scuola e nella vita.

Una delle attività fondamentali dell'informatica è la \textbf{programmazione}: comprendere un problema, progettare una soluzione e scrivere, passo dopo passo, una sequenza di istruzioni che possono essere eseguite da un computer. La parte più pratica e operativa (inserire materialmente queste istruzioni nel computer) è detta dagli informatici \textit{coding}. È una parola molto usata a scuola negli ultimi anni (a volte, con un repentino cambio di significato, addirittura per indicare una nuova disciplina), ma agli informatici non piace, perché riduce l'informatica al solo ``scrivere comandi'', perdendo di vista la comprensione del problema e la progettazione della soluzione. È come paragonare la scrittura creativa all'atto di trascrivere un testo su carta!

Quindi, come sottolineato anche dall'Accademia della Crusca, \textit{coding}, pensiero computazionale e informatica non sono assolutamente sinonimi \parencite{lodi-martini-2024}: l'informatica è una disciplina, mentre il pensiero computazionale è il modo di pensare (e di approcciare i problemi) tipico di chi se ne occupa. La programmazione è un modo con cui gli informatici esprimono le loro soluzioni affinché queste vengano eseguite da un computer. Già ti immaginiamo obiettare: ``Voi la fate facile, ma come pensa un informatico?''  \emoji{thinking-face}

\subsection{Pensiero computazionale }\label{pensiero-computazionale}

Negli anni, sono state proposte innumerevoli definizioni di pensiero computazionale; riportiamo qui quella che ci sembra più completa:


Il \textbf{pensiero computazionale} è il processo mentale coinvolto nella \textbf{formulazione di un problema} e nell'\textbf{esprimere le sue soluzioni} in una forma che può essere \textbf{effettivamente eseguita} da un \textbf{agente di elaborazione delle informazioni} (umano o macchina).

Troppo astratta? Per comprendere meglio di che si tratta, ti consigliamo di guardare il video Exact Instructions Challenge\footnote{\linkbreve{panino}}, di Josh Darnit, su YouTube. Il video è in inglese, ma è facile da comprendere. In alternativa, puoi guardare un vecchio (ma incredibilmente efficace) video del maestro Manzi, in italiano\footnote{\linkbreve{manzi} \textbf{(DAL MINUTO 8:44)}}.

Nel primo video, Josh propone ai suoi figli, Johanna ed Evan, una sfida: scrivere una lista di istruzioni per preparare un sandwich al burro di arachidi e marmellata! Una volta che i figli le avranno scritte, Josh le eseguirà alla lettera. I bambini iniziano con una serie di istruzioni incomplete, che non permettono di raggiungere l'obiettivo, e man mano le perfezionano. Dopo alcuni tentativi, Johanna riesce nell'impresa, mentre Evan, più piccolo, si arrabbia.

Particolarmente interessanti sono alcuni momenti del video (viene indicato il minutaggio del momento del video a cui si fa riferimento):

\begin{itemize}
\item
  Evan ha scritto: ``Prendi una fetta di pane e spalmala in giro con il coltello.'' In queste istruzioni, però, non viene menzionato il burro di arachidi. Quando Josh glielo fa notare, Evan esclama ``Hold on'' (aspetta!), e sul suo volto emerge la consapevolezza che le istruzioni non sono complete. \textbf{1:06}
\item
  Johanna ha scritto: ``Con il coltello, prendi un po' di burro di arachidi,'' e, quando sente il padre leggere la frase ad alta voce, commenta: ``Un po' significa un sacco.'' Alla domanda sorpresa di Josh, ``Un po' significa un sacco?'', lei risponde: ``Nel mio mondo.'' \textbf{1:55}
\item
  Più avanti, Josh legge l'istruzione: ``Inserisci il coltello nel barattolo'' e lo inserisce con la lama rivolta verso l'alto. Lo stupore e il disappunto si leggono sul volto di Johanna: dopo vari tentativi, pensava di aver formulato la soluzione in modo preciso, ma si rende conto che c'era ancora un passaggio ambiguo. \textbf{3:18}
\end{itemize}

Partendo dal video, proviamo a spiegare meglio alcuni aspetti chiave del pensiero computazionale.


\subsubsection{Riflettere sul procedimento di risoluzione dei problemi}

Se il problema (l'obiettivo) dato ai figli di Josh fosse stato quello di preparare un panino al burro di arachidi e marmellata, il risultato sarebbe stato, appunto\ldots{} lo specifico panino da loro preparato! Qui la loro sfida è invece descrivere un modo (una ``ricetta'') per preparare questo tipo di panini. A Johanna ed Evan non basta ottenere un risultato (anzi, non serve: a quello penserà il padre preparando il panino), ma devono pensare al procedimento per ottenerlo.


Questo è un particolare approccio alla soluzione dei problemi tipico dell'informatica: astrarre la soluzione in modo che sia applicabile a tutti i problemi ``di quel tipo'' (Josh potrà seguire la ricetta ogni volta che vuole prepararsi un panino al burro di arachidi e marmellata).


\subsubsection{Esprimere una soluzione}


In questa sfida, descrivere una soluzione (una ricetta) richiede di mettersi nei panni dell'altro. È necessario esprimere la soluzione come una sequenza di istruzioni elementari che possano essere comprese, in modo non ambiguo, dall'esecutore. Quando si descrive un procedimento (o qualsiasi argomento) è fuorviante fare riferimento al ``proprio mondo'', come fa Johanna; bisogna invece avere un'idea chiara delle istruzioni e dei concetti che l'esecutore è in grado di comprendere. Il padre ``fa il finto tonto'', cercando deliberatamente le ambiguità nella descrizione. Non avendo compreso pienamente le regole del gioco, Evan si arrabbia perché sa che l'esecutore è intelligente e conosce già la soluzione; perciò si aspetta che il padre interpreti le sue istruzioni, anche se ambigue.

Il computer è l'esecutore più ``stupido'' di tutti: conosce solo pochissime istruzioni elementari e non aggiunge sue interpretazioni, ma ha il vantaggio di eseguirle sempre allo stesso modo, in modo totalmente prevedibile. Spiegare un procedimento a un computer richiede quindi una comprensione profonda del problema e della soluzione, tenendo sempre conto del livello di comprensione del proprio interlocutore.

Qui il padre ``fa il computer'': esegue alla lettera ciò che legge, senza aggiungere o interpretare.


\subsubsection{Scomporre il problema e le sue soluzioni}


Nel video si nota che Johanna individua subito la necessità di scomporre il problema in due parti: imburrare una fetta e spalmare la marmellata sull'altra, per poi ricomporre le soluzioni parziali nel panino completo. Evan, più giovane, all'inizio fa più fatica; nella sua seconda soluzione, infatti, viene fuori un pasticcio, poiché il problema non è stato scomposto in modo efficace.

È interessante notare che la necessità di scomporre problemi in sottoproblemi non è un'esigenza del computer, ma serve a noi per affrontare una complessità che, se non domata, andrebbe oltre le nostre capacità cognitive.


\subsubsection{Ottenere un feedback dai propri errori}


Ogni volta che la ricetta viene ``eseguita'', Johanna ed Evan ricevono un feedback sulla loro soluzione: osservano che alcune parti funzionano, mentre altre richiedono ulteriori aggiustamenti. La possibilità di testare e migliorare le proprie soluzioni è un aspetto cruciale di questa sfida, poiché consente ai ragazzi di imparare dai propri errori. Senza l'opportunità di imparare dagli errori, infatti, non può esserci un vero apprendimento.

In informatica, questa attività è nota come \textbf{debugging}; ancora una volta, la capacità di analizzare i propri errori è una competenza fondamentale in ogni campo del sapere, dalla matematica alla letteratura. In informatica, tuttavia, il feedback arriva ``prima''. Non serve aspettare che l'insegnante corregga, non c'è in fondo al libro un ermetico numero come risultato. Una volta scritto il programma, si può eseguirlo, ``vedere l'effetto che fa'' e spesso anche osservare qual è l'effetto di ogni singolo passaggio, per verificare se il programma funziona come previsto e, in caso, apportare correzioni.

\subsubsection{Una riflessione conclusiva sul pensiero computazionale}

Ti riproponiamo la definizione di pensiero computazionale:


Il \textbf{pensiero computazionale} è il processo mentale coinvolto nella \textbf{formulazione di un problema} e nell'\textbf{esprimere le sue soluzioni} in una forma che può essere \textbf{effettivamente eseguita} da un \textbf{agente di elaborazione delle informazioni} (umano o macchina).


Questa definizione evidenzia subito un punto cruciale: il pensiero computazionale non riguarda l'esecuzione di procedure, ma si concentra sulla capacità di formulare problemi, progettare soluzioni e rappresentarle in modo chiaro e condivisibile. Si tratta di un particolare modo di risolvere problemi, che potremmo definire \emph{problem solving computazionale}.

Nella scuola, viene dedicato molto tempo (a volte troppo!) al  rafforzamento delle competenze procedurali: imparare tecniche, ripetere esercizi, applicare regole. I bambini e le bambine assumono il ruolo di \textbf{esecutori}. Tutti aspetti importanti, ma che da soli non esauriscono ciò che significa affrontare e risolvere un problema. Il pensiero computazionale invita a fare un passo oltre: non solo ``seguire procedure'', ma anche discutere i problemi, cercare soluzioni generali, analizzare gli errori e trasformarli in occasioni di apprendimento.

Forse già promuovi questo approccio nelle tue attività quotidiane, quando chiedi ad alunni e alunne di spiegare un procedimento, di affrontare un compito complesso a piccoli passi, o di riflettere insieme sugli errori. L'obiettivo di questo testo è proprio aiutarti a riconoscere, in queste pratiche, l'essenza dell'informatica e a rafforzarla sempre di più, per offrire ai bambini e alle bambine competenze che, nel mondo di oggi--fatto di dati, algoritmi e automazione--sono fondamentali per sviluppare un rapporto critico e consapevole con la tecnologia.

\section{Visione didattica }\label{visione-didattica}

\subsection{Informatica senza e con computer }\label{informatica-senza-e-con-computer}

Nei capitoli seguenti, proponiamo attività in due modalità.

\begin{itemize}
\item
  \textbf{Unplugged}: attività \textbf{senza computer}, incentrate sull'uso del corpo, sul linguaggio, sul gioco e sulla collaborazione. Si realizzano con materiali semplici (griglie, carte, pennarelli) e materializzano i concetti dell'informatica in forme concrete e accessibili. Queste attività permettono anche di evitare la visione errata dell'informatica come ``la scienza dei computer''.
\item
  \textbf{Con strumenti digitali}: attività che introducono gradualmente \textbf{ambienti di programmazione} visuale (cioè con istruzioni grafiche, non testuali) e \textbf{robot educativi}. In queste attività, i bambini e le bambine hanno un riferimento tangibile e concreto delle idee sperimentate nelle attività unplugged: concetti astratti come sequenze, condizioni, ripetizioni si materializzano in comportamenti di personaggi che si muovono sullo schermo o di robot che seguono comandi. Queste attività permettono di sperimentare concretamente cosa significhi dare istruzioni a un esecutore automatico ``stupido e intransigente''.
\end{itemize}

Le proposte didattiche seguono il naturale percorso evolutivo degli alunni e delle alunne: dalle esperienze corporee e motorie, dal gioco simbolico e dalla narrazione, fino a raggiungere gradualmente forme sempre più astratte di rappresentazione.

\subsection{Il costruzionismo }\label{il-costruzionismo}

Quello che oggi chiamiamo ``pensiero computazionale'' affonda le sue radici negli studi che Seymour Papert fece già negli anni Settanta, quando i computer iniziarono a entrare nelle scuole. Papert ebbe un ruolo duplice e originale: da matematico e informatico, seppe riconoscere il valore dell'informatica non solo come disciplina da apprendere, ma anche come ambiente privilegiato per pensare, esplorare e costruire conoscenza; da pedagogo, tradusse i principi astratti del costruttivismo di Jean Piaget in pratiche educative concrete, capaci di dare forma all'apprendimento quotidiano degli alunni e alunne.

Da questo intreccio nasce il \textbf{costruzionismo}, la sua teoria dell'apprendimento. Come il costruttivismo, anche il costruzionismo sostiene che la conoscenza non si riceve passivamente, ma si costruisce attivamente attraverso l'esperienza. La novità introdotta da Papert è l'idea che l'apprendimento diventi particolarmente efficace quando prende la forma di un'attività creativa orientata alla produzione di artefatti significativi: prodotti che hanno valore per chi li realizza o per la comunità di cui fanno parte.

Papert chiama questi artefatti ``oggetti con cui pensare'': strumenti che rendono le idee visibili, manipolabili, discutibili e migliorabili, facilitando la costruzione di modelli mentali. Possono essere concreti o astratti--un disegno, un racconto, una costruzione con i mattoncini, una sequenza di istruzioni, un programma al computer--ma ciò che conta è che offrano a chi apprende un terreno fertile su cui sperimentare, riflettere e crescere.


\subsubsection{Costruzionismo e informatica}


Nella prefazione del suo libro più famoso, Mindstorms, \textcite{papert-mindstorms-1980} afferma che da bambino era affascinato dagli ingranaggi--smontava qualunque oggetto meccanico per capirne il funzionamento. Il suo modello mentale sugli ingranaggi gli ha permesso di comprendere concetti matematici più complessi, come le proporzioni, le frazioni, fino alle equazioni in più variabili:


\emph{``Credo che l'aver tanto giocato con i differenziali sia stato più efficace per la mia comprensione della matematica, di tutto quello che mi è stato insegnato alla scuola elementare. Gli ingranaggi, servendomi da modelli, hanno fatto entrare nella mia mente idee che altrimenti sarebbero restate astratte.''}


Nella scelta di questi ``oggetti con cui pensare'', però, non c'è solo un aspetto conoscitivo: è sempre in gioco una fondamentale componente affettiva. La passione di Papert per gli ingranaggi è personale e non riproducibile.


\emph{``Una Maria Montessori dei nostri giorni potrebbe proporre, se la mia storia l'avesse convinta, di creare una confezione di ingranaggi per bambini. Così ogni bambino potrebbe fare l'esperienza che ho fatto io. Ma sperarlo sarebbe trascurare l'essenza della storia: io mi ero innamorato degli ingranaggi. È dunque qualcosa che non può essere ridotto in termini soltanto cognitivi. Era avvenuto qualcosa di molto personale, e non si può presumere che ciò si ripeterebbe per altri bambini esattamente nella stessa forma.''}


Il computer, per Papert, diventa il più flessibile fra questi ``oggetti con cui pensare'', una sorta di oggetto universale, capace di rappresentare e simulare qualsiasi sistema o concetto. Come una volta gli ingranaggi hanno offerto a lui stesso una comprensione intuitiva di concetti complessi, ora l'informatica può fornire a ogni studente un insieme flessibile di strumenti con cui esplorare e comprendere idee astratte. Il computer non è quindi solo un mezzo per eseguire calcoli o simulazioni, ma uno spazio di esplorazione in cui gli alunni e le alunne possono creare e modificare in prima persona.


\emph{``La mia tesi potrebbe essere sintetizzata nel modo seguente: quello che gli ingranaggi non possono fare, il computer lo potrebbe. Il computer è il Proteo delle macchine. La sua essenza è la sua universalità, il suo potere di simulare. Poiché può assumere migliaia di forme e può essere utilizzato per migliaia di funzioni, può soddisfare migliaia di gusti.''}


Il computer ha anche un altro ruolo nell'apprendimento: consente agli alunni e alunne di vedere immediatamente il risultato delle proprie azioni--ogni comando o modifica genera una risposta visibile e immediata, che stimola ulteriori tentativi e incoraggia l'auto-correzione. Questo ciclo di azione-feedback-miglioramento rende l'apprendimento più efficace, consentendo agli alunni e alle alunne di sperimentare. In questo senso, il feedback immediato del computer diventa un potente ``oggetto con cui pensare'', uno strumento che aiuta a interiorizzare e applicare concetti in un processo di scoperta continua. Questo aspetto rende l'uso del computer particolarmente adatto alla pedagogia costruzionista, favorendo un apprendimento in cui il discente può iterare, riflettere e progredire autonomamente nel proprio percorso di scoperta.

Come si collega tutto questo alle discipline studiate nel primo ciclo? Papert utilizza il termine ``Mathland'' per indicare un paese immaginario in cui la matematica si impara con la stessa naturalezza con cui si impara a parlare, vivendo immersi in una lingua. L'informatica può diventare questo ambiente speciale: non un insieme di esercizi astratti, ma un luogo in cui i concetti matematici vengono costruiti tramite l'interazione con artefatti significativi. Per esempio, quando i bambini e le bambine programmano un piccolo robot, esplorano la geometria muovendosi sul piano e facendo uso di nozioni come angoli (a tale proposito, si vedano le attività nelle guide di classe terza primaria del progetto PerContare \urloriginale{https://www.percontare.it/}). In questo modo, la matematica non resta ``fuori'' dalle loro esperienze, ma diventa parte viva dei loro giochi, dei loro progetti e delle loro invenzioni.

Papert sottolinea che lo stesso approccio può estendersi anche al linguaggio. Così come la matematica può diventare parte viva dei giochi e delle esplorazioni, l'informatica può sostenere le capacità narrative ed espressive. Costruire storie interattive, far dialogare personaggi o programmare piccoli robot che recitano un racconto significa unire informatica e storytelling: gli alunni e le alunne imparano a dare ordine agli eventi, a strutturare sequenze, a esprimere le proprie idee in forma chiara e condivisibile.


\subsubsection{Costruzionismo e apprendimento}


Nel costruzionismo di Papert, l'obiettivo è ribaltare il processo tradizionale di insegnamento, passando da un approccio in cui il docente trasferisce conoscenze a uno in cui studenti e studentesse diventano protagonisti del proprio apprendimento. Invece di fornire informazioni già pronte, l'insegnante ha il compito di creare condizioni che favoriscano la scoperta. Per il costruzionismo, quindi, l'apprendimento è un autentico processo esplorativo, in cui il ruolo dell'insegnante è quello di offrire agli alunni e alunne esperienze e strumenti che consentano loro di approfondire concetti e costruire attivamente la propria conoscenza. Tradurre questa prospettiva in pratica significa organizzare attività che non siano solo esercizi da svolgere, ma vere esperienze di costruzione e riflessione. Alcuni principi chiave:

\begin{itemize}
\item
  \textbf{Ribaltare il processo di apprendimento}. Bisogna porsi sempre la domanda: \textit{come posso fare un passo indietro e rendere gli alunni e le alunne protagonisti?} L'obiettivo non è fornire risposte, ma stimolare la ricerca e la costruzione di soluzioni personali.

\item
  \textbf{Personalizzazione e motivazione}. Le attività hanno più valore se lasciano spazio agli interessi personali. Per esempio, programmare con Scratch un biglietto di auguri per la propria migliore amica non è solo un modo per imparare comandi, strutture di ripetizione o variabili, ma è anche (soprattutto) un modo per esprimere qualcosa di personale e rilevante. Questo rende la costruzione significativa e motivante, proprio come per Papert lo erano stati gli ingranaggi.

\item
  \textbf{Valorizzare l'errore come motore creativo}. In un approccio costruzionista l'errore non è un fallimento, ma una scoperta. Ogni volta che un programma non funziona o che il risultato non è quello atteso, si apre uno spazio di ricerca: cos'è successo? Cosa posso cambiare? Spesso gli errori portano a soluzioni nuove e inattese. In questo senso, l'errore diventa un processo naturale di esplorazione e non un motivo di frustrazione.

\item
  \textbf{Iterazione e miglioramento continuo}. Non esiste ``la soluzione giusta e basta'': ogni costruzione può essere raffinata, estesa, modificata. Ogni processo può essere migliorato. Abituare gli alunni e le alunne a ripensare al proprio lavoro, cercando di renderlo più chiaro o più efficace, rafforza l'idea che il sapere è un processo in evoluzione.

\item
  \textbf{Discussione e riflessione collettiva}. Il momento finale dell'attività è spesso il più importante: non fermarsi al ``compito svolto'', ma dedicare tempo a discutere insieme. Che cosa ha funzionato? Dove ci siamo confusi? Come avremmo potuto fare meglio? Questo passaggio trasforma l'esperienza in apprendimento riflessivo: non solo ``abbiamo fatto'', ma ``abbiamo capito cosa significa fare''.

\item
  \textbf{Artefatti visibili e condivisi}. Ogni attività dovrebbe concludersi con un prodotto tangibile (un disegno, una sequenza scritta, una mappa, un programma). Questo oggetto diventa la traccia del ragionamento fatto, può essere mostrato, discusso, confrontato. Così idee, scelte e strategie diventano visibili e condivisibili.

\item
  \textbf{Uso creativo degli strumenti digitali}. Non proporre gli strumenti solo per eseguire esercizi guidati, ma lasciare che diventino mezzi espressivi: con Scratch si può raccontare una fiaba interattiva, con un robot mettere in scena una coreografia, con un editor di testi costruire un giornalino di classe.

\item
  \textbf{Ruolo dell'insegnante come facilitatore}. Nel costruzionismo, l'insegnante non fornisce soluzioni già pronte, ma crea condizioni per la scoperta. In pratica: non corregge subito un errore, ma pone domande (``Cosa ti aspettavi che accadesse?''), rilancia le idee degli alunni e delle alunne, favorisce il confronto tra pari. È un ruolo più simile a quello di un allenatore che stimola e incoraggia, più che a quello di un esaminatore che giudica.
\end{itemize}


\section{Suggerimenti per ulteriori attività}\label{suggerimenti-attivita-cap1}

\printbibliography[heading=none,keyword={cap1},category={attivita}]

\section{Riferimenti}\label{riferimenti}

\printbibliography[heading=none,keyword={cap1},notcategory={attivita}]
